% !TEX program = xelatex
\documentclass[aspectratio=169,12pt]{ctexbeamer}

% ── Theme ────────────────────────────────────────────────────
\usetheme{Madrid}
\usecolortheme{whale}
\setbeamertemplate{navigation symbols}{}
\setbeamertemplate{footline}[frame number]
\setbeamertemplate{caption}[numbered]

% ── Packages ──────────────────────────────────────────────────
\usepackage{graphicx}
\usepackage{booktabs}
\usepackage{amsmath}
\usepackage{xcolor}
\usepackage{hyperref}

\graphicspath{{../outputs/}}

% ── Title ──────────────────────────────────────────────────────
\title{基于股吧散户情绪的\\另类 Alpha 因子策略}
\subtitle{网络爬虫 · 情感分析 · 回归建模}
\author{XXX}
\institute{上海财经大学 · 统计与管理学院}
\date{\today}

\begin{document}

% ═══════════════════ SLIDE 1 — Title ══════════════════════════
\begin{frame}
    \titlepage
\end{frame}

% ═══════════════════ SLIDE 2 — 目录 ══════════════════════════
\begin{frame}{目录}
    \tableofcontents
\end{frame}

% ═══════════════════ SLIDE 3 — 背景 ══════════════════════════
\section{研究背景}

\begin{frame}{研究背景}
    \begin{columns}[T]
        \column{0.55\textwidth}
        \textbf{量化投资的演变}
        \begin{itemize}
            \item 传统 Alpha 因子：价量、财务、分析师一致预期
            \item 另类数据（Alternative Data）：社交媒体、卫星图像、信用卡交易
            \item 中国A股特色：散户占比高，情绪驱动性强
        \end{itemize}
        \vspace{0.5cm}
        \textbf{为什么选择股吧？}
        \begin{itemize}
            \item 东方财富股吧：日均数万条新帖
            \item 淘股吧：深度讨论，专业散户聚集地
            \item 文本情绪 → 短期股价预测？
        \end{itemize}
        \column{0.4\textwidth}
        \centering
        \includegraphics[width=\textwidth]{platform_comparison.png}
    \end{columns}
\end{frame}

% ═══════════════════ SLIDE 4 — 研究问题 ══════════════════════
\section{研究问题}

\begin{frame}{研究问题}
    \begin{block}{核心问题}
        \begin{enumerate}
            \item 股吧散户情绪能否显著预测个股未来一日收益？
            \item 不同情绪因子（均值 vs 分歧 vs 热度）预测力有何差异？
            \item Ridge / LASSO / LightGBM 在该任务上表现如何？
        \end{enumerate}
    \end{block}
    \vspace{0.5cm}
    \begin{alertblock}{目标}
        构建从爬虫到建模到回测的完整端到端量化研究流水线。
    \end{alertblock}
\end{frame}

% ═══════════════════ SLIDE 5 — 方法总览 ══════════════════════
\section{方法论总览}

\begin{frame}{技术路线图}
    \centering
    \begin{enumerate}[<+->]
        \item \textbf{数据采集}：Python 爬虫 → PostgreSQL 存储
        \item \textbf{NLP 处理}：FinBERT-Chinese → 情感得分 $[-1,1]$
        \item \textbf{因子构造}：4个日频情绪因子 + 控制变量
        \item \textbf{回归建模}：Ridge / LASSO / LightGBM 对比
        \item \textbf{回测评估}：十分组多空组合 → 夏普比率
    \end{enumerate}
    \vspace{0.5cm}
    \visible<6->{
    \begin{block}{技术栈}
        Python · PostgreSQL · HuggingFace · scikit-learn · LightGBM · Matplotlib
    \end{block}
    }
\end{frame}

% ═══════════════════ SLIDE 6 — 数据采集 ══════════════════════
\section{数据采集}

\begin{frame}{数据采集：网络爬虫}
    \begin{columns}[T]
        \column{0.5\textwidth}
        \textbf{数据源}
        \begin{itemize}
            \item 东方财富股吧（主流，高流量）
            \item 淘股吧（深度讨论）
            \item baostock（行情数据）
        \end{itemize}
        \vspace{0.3cm}
        \textbf{技术要点}
        \begin{itemize}
            \item requests + BeautifulSoup
            \item GBK/UTF-8 自动检测
            \item 速率限制：2--3s/页
            \item 无需 Cookie（东方财富）
        \end{itemize}
        \column{0.45\textwidth}
        \centering
        \includegraphics[width=\textwidth]{posts_over_time.png}
    \end{columns}
\end{frame}

% ═══════════════════ SLIDE 7 — NLP ══════════════════════════
\section{NLP 情感分析}

\begin{frame}{NLP 情感分析：FinBERT-Chinese}
    \begin{columns}[T]
        \column{0.55\textwidth}
        \textbf{模型选择}
        \begin{itemize}
            \item \texttt{yiyanghkust/finbert-tone-chinese}
            \item BERT-base 架构，约 4.1 亿参数
            \item 中文金融分析师报告微调
            \item 测试准确率 88\%
        \end{itemize}
        \vspace{0.3cm}
        \textbf{情感得分映射}
        \[
        \text{score} = P(\text{positive}) - P(\text{negative})
        \]
        范围 $[-1, +1]$
        \column{0.4\textwidth}
        \centering
        \includegraphics[width=\textwidth]{sentiment_distribution.png}
    \end{columns}
\end{frame}

% ═══════════════════ SLIDE 8 — 因子工程 ══════════════════════
\section{因子工程}

\begin{frame}{因子工程：4个日频情绪因子}
    \begin{columns}[T]
        \column{0.55\textwidth}
        \begin{enumerate}[F1]
            \item \textbf{关注度异常}
            \[
            \frac{\text{日发帖量}}{\text{20日均量}}
            \]
            \item \textbf{情绪均值}
            \[
            \frac{1}{N}\sum \text{sentiment}
            \]
            \item \textbf{情绪分歧}
            \[
            \text{std}(\text{sentiment})
            \]
            \item \textbf{互动热度}
            \[
            \text{avg}(\text{回复数})
            \]
        \end{enumerate}
        \column{0.4\textwidth}
        \textbf{控制变量}
        \begin{itemize}
            \item 当日收益率
            \item 20日波动率
            \item 成交量异常
            \item 换手率
            \item 市场因子（沪深300）
        \end{itemize}
        \vspace{0.3cm}
        \textbf{目标变量}
        \begin{itemize}
            \item 次日收益率
        \end{itemize}
    \end{columns}
\end{frame}

\begin{frame}{因子相关性分析}
    \centering
    \includegraphics[width=0.55\textwidth]{factor_correlation.png}
\end{frame}

% ═══════════════════ SLIDE 9 — 建模 ══════════════════════════
\section{回归建模}

\begin{frame}{回归建模：三种模型对比}
    \begin{columns}[T]
        \column{0.33\textwidth}
        \textbf{Ridge 回归}
        \begin{itemize}
            \item $L_2$ 正则化
            \item 5 折交叉验证
            \item $\alpha$ 网格搜索
            \item 压缩但不消除特征
        \end{itemize}
        \column{0.33\textwidth}
        \textbf{LASSO 回归}
        \begin{itemize}
            \item $L_1$ 正则化
            \item 自动特征选择
            \item LassoCV 自动调参
            \item 稀疏解
        \end{itemize}
        \column{0.33\textwidth}
        \textbf{LightGBM}
        \begin{itemize}
            \item 梯度提升树
            \item 非线性建模
            \item Early stopping
            \item 500 棵树
        \end{itemize}
    \end{columns}
    \vspace{0.5cm}
    \begin{block}{数据划分}
        时间序列划分（前 80\% 训练，后 20\% 测试）· 标准化处理
    \end{block}
\end{frame}

\begin{frame}{模型性能对比}
    \centering
    \includegraphics[width=0.65\textwidth]{feature_importance.png}
    \vspace{0.3cm}

    \begin{table}
        \centering
        \begin{tabular}{lrrr}
            \toprule
            \textbf{模型} & $\boldsymbol{R^2}$ & \textbf{MSE} & \textbf{MAE} \\
            \midrule
            Ridge & \texttt{RIDGE\_R2} & \texttt{RIDGE\_MSE} & \texttt{RIDGE\_MAE} \\
            LASSO & \texttt{LASSO\_R2} & \texttt{LASSO\_MSE} & \texttt{LASSO\_MAE} \\
            LightGBM & \texttt{LGB\_R2} & \texttt{LGB\_MSE} & \texttt{LGB\_MAE} \\
            \bottomrule
        \end{tabular}
    \end{table}
\end{frame}

% ═══════════════════ SLIDE 10 — 回测 ══════════════════════════
\section{回测评估}

\begin{frame}{回测策略：十分组多空组合}
    \begin{columns}[T]
        \column{0.5\textwidth}
        \textbf{策略构建}
        \begin{enumerate}
            \item 每日用模型预测所有股票收益
            \item 按预测值排序 → 等分 10 组
            \item Group 1：做多（高预测收益）
            \item Group 10：做空（低预测收益）
            \item 等权配置
        \end{enumerate}
        \column{0.45\textwidth}
        \textbf{评估指标}
        \begin{itemize}
            \item 年化收益率
            \item 夏普比率
            \item 最大回撤
            \item 胜率
            \item 分组收益单调性
        \end{itemize}
    \end{columns}
\end{frame}

\begin{frame}{回测结果：累计收益曲线}
    \centering
    \includegraphics[width=0.75\textwidth]{backtest_cumulative.png}
\end{frame}

\begin{frame}{回测结果：十分组收益}
    \centering
    \includegraphics[width=0.7\textwidth]{backtest_decile_returns.png}

    \vspace{0.3cm}
    \begin{table}
        \centering
        \small
        \begin{tabular}{lrrrr}
            \toprule
            \textbf{模型} & \textbf{年化收益} & \textbf{夏普比率} & \textbf{最大回撤} & \textbf{胜率} \\
            \midrule
            Ridge & \texttt{R\_RET} & \texttt{R\_SR} & \texttt{R\_MDD} & \texttt{R\_WR} \\
            LASSO & \texttt{L\_RET} & \texttt{L\_SR} & \texttt{L\_MDD} & \texttt{L\_WR} \\
            \bottomrule
        \end{tabular}
    \end{table}
\end{frame}

% ═══════════════════ SLIDE 11 — 结论 ══════════════════════════
\section{结论}

\begin{frame}{主要结论}
    \begin{block}{发现}
        \begin{enumerate}
            \item \textbf{技术可行}：端到端流水线可在个人计算机运行
            \item \textbf{情感因子有区分度}：截面选股具备一定能力，Long-Short 组合可获超额收益
            \item \textbf{Ridge 最稳健}：在该任务上略优于 LASSO 和 LightGBM
        \end{enumerate}
    \end{block}
    \vspace{0.3cm}
    \begin{alertblock}{核心贡献}
        构建了从数据获取到模型评估的完整可复现量化研究框架
    \end{alertblock}
\end{frame}

% ═══════════════════ SLIDE 12 — 局限 ══════════════════════════
\section{局限与展望}

\begin{frame}{局限性与未来方向}
    \begin{columns}[T]
        \column{0.48\textwidth}
        \textbf{当前局限}
        \begin{itemize}
            \item 样本量受限（一年数据）
            \item FinBERT 领域差异（分析师 vs 散户口语）
            \item 仅日频因子
            \item 未计入交易成本
        \end{itemize}
        \column{0.48\textwidth}
        \textbf{未来方向}
        \begin{itemize}
            \item 扩展至创业板、科创板
            \item Fine-tune 股吧领域情感模型
            \item 加入盘中高频情绪因子
            \item 考虑交易成本的净收益评估
            \item 引入深度学习时序模型（LSTM/Transformer）
        \end{itemize}
    \end{columns}
\end{frame}

% ═══════════════════ SLIDE 13 — 感谢 ══════════════════════════
\begin{frame}[plain]
    \centering
    \vspace{2cm}
    {\Huge 谢谢！}\\[1cm]
    {\large 欢迎提问}\\[2cm]
    {\small 项目代码：github.com/.../retail-sentiment-alpha}
\end{frame}

\end{document}
